\section*{Chapitre \ref{ch:b3:hydrogeology} --- Hydrogéologie}

\begin{solution}{exo:b3:hydrogeology:1}
Premier capteur~: $z = 120 - 40 = \qty{80}{m}$, $h = 120 - 12.5 =
\qty{107.5}{m}$, charge de pression $107.5 - 80 = \qty{27.5}{m}$. Deuxième
capteur~: $z = \qty{40}{m}$, $h = \qty{105.0}{m}$, charge de pression
\qty{65.0}{m}. La charge est plus élevée au capteur peu profond, donc l'eau
s'écoule vers le bas, sous un gradient $(107.5 - 105.0)/40 = \num{0.0625}$~:
les piézomètres se trouvent dans une zone d'alimentation. La pression seule aurait
suggéré le contraire.
\end{solution}

\begin{solution}{exo:b3:hydrogeology:2}
$q = K\Delta h/L = 12\times0.30/0.50 = \qty{7.2}{m/d}$. La section vaut
$\pi\times0.05^2 = \qty{7.85e-3}{m^2}$, donc $Q = \qty{0.0565}{m^3/d}$, soit
$56.5/24 = \qty{2.36}{L/h}$. La \omterm{def:b3:hydrogeology:seepage}{vitesse effective} vaut $7.2/0.35 =
\qty{20.6}{m/d}$.
\end{solution}

\begin{solution}{exo:b3:hydrogeology:3}
$K = 2.5/86400 = \qty{2.894e-5}{m/s}$, et $k = K\mu/(\rho g) =
\num{2.894e-5}\times\num{1.306e-3}/(999.7\times9.81) =
\qty{3.85e-12}{m^2}$. À \qty{20}{\celsius}, $K$ est multiplié par
$(998.2/\num{1.002e-3})/(999.7/\num{1.306e-3}) = \num{1.30}$~: $K =
\qty{3.25}{m/d}$. La variation tient presque entièrement à la viscosité.
\end{solution}

\begin{solution}{exo:b3:hydrogeology:4}
$T = 2.5\times180 = \qty{450}{m^2/d}$. La dilatation de l'eau donne
$\num{9e-7}\times180 = \num{1.62e-4}$, soit environ \qty{80}{\percent} de $S =
\num{2e-4}$~; la compression du calcaire fournit le reste.
\end{solution}

\begin{solution}{exo:b3:hydrogeology:5}
D'après la \cref{prop:b3:hydrogeology:thiem}, $T = Q\ln(r_2/r_1)/(2\pi(s_1 -
s_2)) = 1500\times\ln10/(2\pi\times1.53) = \qty{359}{m^2/d}$, et $K =
359/40 = \qty{9.0}{m/d}$.
\end{solution}

\begin{solution}{exo:b3:hydrogeology:6}
$Q/4\pi T = 1200/(4\pi\times300) = \qty{0.318}{m}$. Après 1 jour, $u =
300^2\times\num{5e-4}/(4\times300\times1) = \num{0.0375}$ et $W = -0.5772
+ 3.2834 + 0.0375 - 0.0004 = \num{2.743}$~: $s = \qty{0.873}{m}$. Après 10
jours, $u = \num{0.00375}$, $W = -0.5772 + 5.5860 + 0.0038 = \num{5.013}$~:
$s = \qty{1.596}{m}$. Un temps dix fois plus long ne donne pas deux fois
le rabattement.
\end{solution}

\begin{solution}{exo:b3:hydrogeology:7}
$u \le 0.05$ lorsque $t \ge r^2S/(4T\times0.05)$. À \qty{60}{m}~: $3600\times
\num{2e-4}/90 = \qty{0.008}{d}$, soit \qty{11.5}{min}. À \qty{500}{m}~:
\qty{0.556}{d}, \qty{800}{min}, soit plus de 13 heures. Le temps croît comme $r^2$~:
dans un essai d'un ou deux jours, les puits d'observation utilisés pour la droite
doivent se trouver à quelques centaines de mètres du puits pompé.
\end{solution}

\begin{solution}{exo:b3:hydrogeology:8}
À mi-distance, chaque puits est à \qty{200}{m}~: $u = 200^2\times\num{2e-4}/(4
\times450\times10) = \num{4.44e-4}$, $W = \num{7.142}$, $s = 0.382\times
7.142 = \qty{2.728}{m}$ de chaque puits, soit \qty{5.456}{m} au total. À
\qty{400}{m}, $u = \num{1.78e-3}$, $W = \num{5.757}$, $s =
\qty{2.199}{m}$. Chaque puits abaisse l'eau de l'autre de
\qty{2.2}{m}~: sa pompe doit remonter l'eau de \qty{2.2}{m} de plus, ce qui
coûte de l'énergie, et le rabattement qui subsiste avant que la crépine ne
se découvre est plus faible.
\end{solution}

\begin{solution}{exo:b3:hydrogeology:9}
Quatre lignes de courant entre la base du barrage et le substratum délimitent $N_f = 5$
chenaux~; neuf équipotentielles divisent les \qty{10}{m} en $N_d = 10$
chutes. $q = KHN_f/N_d = 5\times10\times5/10 = \qty{25}{m^3/d}$ par mètre
de barrage, et \qty{7500}{m^3/d} sous \qty{300}{m}. Le facteur de forme calculé,
0.512 au lieu de 0.5, donnerait \qty{25.6}{m^3/d} par mètre.
\end{solution}

\begin{solution}{exo:b3:hydrogeology:10}
$k = na^2/8 = 0.30\times(\num{5e-5})^2/8 = \qty{9.4e-11}{m^2}$, et $K =
k\rho g/\mu = \num{9.375e-11}\times998.2\times9.81/\num{1.002e-3} =
\qty{9.2e-4}{m/s}$, soit \qty{79}{m/d}, dans la gamme du sable propre, de 0.3 à
\qty{150}{m/d}. Le modèle néglige la tortuosité des trajets, plus longs que
les tubes droits (elle abaisse $K$), les gorges entre les pores, plus étroites
que les pores (effet à la baisse), les pores en cul-de-sac qui ne transportent
aucun débit (à la baisse), et la dispersion des tailles de pores~: comme un tube
transporte un débit proportionnel à $a^4$, les trajets les plus larges
transportent l'essentiel de l'eau, et une dispersion élève $K$ au-dessus de celui de tubes de taille moyenne.
\end{solution}

\begin{solution}{exo:b3:hydrogeology:11}
(a)~Le logarithme s'annule lorsque $2.25\,Tt = r^2S$, soit en $r_0 =
\sqrt{2.25\,Tt/S}$. (b)~Nappe captive puis nappe libre~:
\[
r_0 = \sqrt{\frac{2.25\times450\times30}{\num{2e-4}}} = \qty{12.3}{km},
\qquad r_0 = \sqrt{\frac{2.25\times450\times30}{0.2}} = \qty{390}{m},
\]
soit $\sqrt{1000} = 31.6$ fois plus petit. (c)~La pression
se propage avec la diffusivité $T/S$, très grande lorsque $S$ est petit~: pendant
l'essai, le cône atteint les failles, les affleurements de l'aquifère et des rivières,
et le rabattement quitte la droite, ce qui indique à
l'hydrogéologue qu'une limite a été atteinte.
\end{solution}

\begin{solution}{exo:b3:hydrogeology:12}
(a)~Vu depuis le référentiel qui se déplace avec l'eau, le front s'étale avec
un écart type $\sigma_x = \sqrt{2Dt} = \sqrt{2\alpha_Lvt}$. Il passe en
$x$ vers $t = x/v$, où $\sigma_x = \sqrt{2\alpha_Lx}$~; à la vitesse
$v$, un étalement $\sigma_x$ en distance correspond à un étalement $\sigma_t = \sigma_x/v$
en temps. (b)~Les temps à \qty{16}{\percent} et à \qty{84}{\percent} encadrent
la moyenne à un écart type de chaque côté~: $\sigma_t = (247.2 -
152.6)/2 = \num{47.3}~\text{jours}$. Alors $\alpha_L = v^2\sigma_t^2/(2x) =
0.25\times47.3^2/200 = \qty{2.8}{m}$. (La courbe a été calculée avec
\qty{3}{m}~: l'approximation normale la sous-estime légèrement.) (c)~$R =
1 + 1.9\times0.32/0.30 = \num{3.03}$~: le soluté se déplace au tiers de la
vitesse de l'eau.
\end{solution}

\begin{solution}{pb:b3:hydrogeology:1}
\textbf{1.} Schiste $363.2 - 282.5 = \qty{80.7}{m}$~; calcaire $282.5 -
102.5 = \qty{180}{m}$.

\textbf{2.} $h = 363.2 - 45.0 = \qty{318.2}{m}$~; charge de pression au sommet
$318.2 - 282.5 = \qty{35.7}{m}$. L'eau s'élève de \qty{35.7}{m} au-dessus du
sommet du calcaire, retenue par le schiste~: l'aquifère est à nappe captive.

\textbf{3.} À partir de $h_2 - h_1 = 550a - 100b = -1.62$ et $h_3 - h_1 = 250a
- 600b = -1.44$~: $a = \num{-2.715e-3}$ et $b = \num{1.269e-3}$ par
mètre.

\textbf{4.} $|\nabla h| = \sqrt{a^2 + b^2} = \num{0.0030}$. L'eau
s'écoule selon $-\nabla h = (\num{2.715e-3}, \num{-1.269e-3})$, vers l'est et
au sud~: azimut $\ang{180} - \arctan(2.715/1.269) = \ang{115}$.

\textbf{5.} Les points situés avant environ \qty{10}{min} se trouvent au-dessus de la droite~: là
$u > 0.05$ et la courbe de Theis se situe au-dessus de son approximation logarithmique
($W > -\gamma - \ln u$, comme le montre la figure de la fonction de puits).

\textbf{6.} La meilleure droite passant par les points à partir de \qty{15}{min} passe
\qty{1.75}{m} à \qty{100}{min} et \qty{2.63}{m} à \qty{1000}{min}~:
$\Delta s = \qty{0.87}{m}$ par cycle logarithmique.

\textbf{7.} $T = 2.3\times2160/(4\pi\times0.872) = \qty{453}{m^2/d}$.

\textbf{8.} La droite descend de \qty{0.87}{m} par cycle à partir de \qty{1.75}{m} à
\qty{100}{min}~: le zéro est atteint $1.75/0.872 = 2.01$ cycles plus tôt,
$t_0 = \qty{0.98}{min} = \qty{6.83e-4}{d}$.

\textbf{9.} $S = 2.25\,Tt_0/r^2 = 2.25\times453\times\num{6.83e-4}/3600 =
\num{1.9e-4}$, dans la gamme des \omterm{def:b3:hydrogeology:aquifer-types}{aquifères à nappe captive}.

\textbf{10.} À \qty{15}{min}, $u = 3600\times\num{1.9e-4}/(4\times453
\times0.0104) = \num{0.036}$, et moins par la suite.

\textbf{11.} $K = 453/180 = \qty{2.5}{m/d}$, un calcaire fracturé.

\textbf{12.} $u = 3600\times\num{2e-4}/(4\times450\times2) =
\num{2.0e-4}$, $W = -0.5772 + 8.5172 + 0.0002 = \num{7.940}$, $s =
0.382\times7.940 = \qty{3.03}{m}$, comme relevé.

\textbf{13.} $u = 10^6\times\num{2e-4}/(4\times450\times30) =
\num{0.0037}$, $W = \num{5.02}$, $s = \qty{1.92}{m}$.

\textbf{14.} $r_0 = \sqrt{2.25\times450\times30/(\num{2e-4})} =
\qty{12.3}{km}$, plus que la largeur de la région. Le \omterm{def:b3:hydrogeology:cone}{cône de rabattement} a atteint
les bords du calcaire, ses affleurements et la faille, qui pourraient le mettre
au contact de roches moins perméables~: l'aquifère infini de Theis ne
le décrit plus, et seul l'enregistrement des rabattements dira si une
limite alimente le puits ou arrête l'écoulement.

\textbf{15.} $2160\times30 = \qty{64800}{m^3}$. Ce volume provient de l'emmagasinement~:
l'eau libérée par la dilatation de l'eau et la compression
du calcaire, sur toute la surface du cône. En intégrant $S\,s$ sur
le cône restitue $Qt$.

\textbf{16.} $u = 10^6\times0.18/(4\times450\times30) = \num{3.33}$, $W
= \num{0.0086}$, $s = \qty{3}{mm}$~: la ferme ne s'en apercevrait pas. Le
pompage puiserait dans les pores près du puits au lieu de la pression
de l'ensemble de l'aquifère.

\textbf{17.} $q = 2.5\times0.0030 = \qty{0.0075}{m/d}$; $v = 0.0075/0.18
= \qty{0.042}{m/d}$, \qty{15}{m} par an.

\textbf{18.} $200/0.0417 = \num{4800}\text{ jours}$, \num{13.1}~ans~; le
polluant, \num{32.9}~ans.

\textbf{19.} À travers un cylindre de rayon $r$, le débit $Q$ traverse la
surface $2\pi rb$ dont les pores représentent la fraction $n_e$~: $\mathrm dr/\mathrm
dt = -Q/(2\pi rbn_e)$, et en intégrant $2\pi bn_er\,\mathrm dr = -Q\,
\mathrm dt$ de $r$ à 0 donne $t = \pi n_ebr^2/Q$. Ici $\pi\times0.18
\times180\times200^2/2160 = \num{1885}~\text{jours}$, \num{5.2}\text{ ans} pour
l'eau et $2.5\times5.2 = \num{12.9}\text{ ans}$ pour le polluant~: le puits
l'amène plus de deux fois plus tôt.

\textbf{20.} $\alpha_L = \qty{20}{m}$ et $\sqrt{2\times20\times200} =
\qty{89}{m}$. La concentration atteint \qty{16}{\percent} de sa
valeur finale \qty{89}{m} en avance sur le front moyen~: ce niveau atteint le
puits lorsque le front moyen se trouve encore à \qty{89}{m}, au bout de $(111/200)^2\times12.9 = \num{4.0}\text{ ans}$, soit un tiers
du temps moyen. Les premières traces comptent pour l'eau potable, non la
moyenne.

\textbf{21.} Pour~: le puits capte le polluant, qui ne peut donc
s'étendre plus loin vers l'aval, et l'eau peut être traitée là où elle est pompée.
Contre~: le pompage attire le polluant vers un puits d'alimentation plus vite
qu'il ne le ferait autrement, et l'eau doit être surveillée et
traitée pendant des décennies.

\textbf{22.} La \omterm{def:b3:hydrogeology:storage}{transmissivité} du calcaire de Varden est $T \approx
\textbf{\qty[text-series-to-math=true]{450}{m^2/d}}$ (\qty{453}{m^2/d}
issue de l'essai).
\end{solution}
